\section{Geometric complexity and orbit closures}
\label{sec:gct}

Let \(V=\operatorname{Mat}_n(\mathbb C)\), put
\(U=\operatorname{Sym}^n(V^*)\), and let \(G=\operatorname{GL}(V)\)
act on forms by \((g\cdot f)(x)=f(g^{-1}x)\). For \(m<n\), embed
\(\operatorname{per}_m\) using \(m^2\) coordinates and choose a further
coordinate \(z\). The two affine orbit closures are
\begin{equation*}
  D_n=\overline{G\cdot\det_n},\qquad
  P_{m,n}=\overline{G\cdot(z^{n-m}\operatorname{per}_m)}.
\end{equation*}
A coefficient polynomial that vanishes on the determinant orbit and
evaluates to a nonzero value on the padded permanent certifies
\(P_{m,n}\nsubseteq D_n\).

Both closures are cones. If \(P_{m,n}\subseteq D_n\), the homogeneous
degree-\(d\) pieces of their
coordinate rings admit a surjective \(G\)-equivariant map
\begin{equation*}
  \mathbb C[D_n]_d \longrightarrow \mathbb C[P_{m,n}]_d.
\end{equation*}
These are finite-dimensional rational representations of a reductive
group over \(\mathbb C\). Complete reducibility gives
\(m_{\lambda,d}(P_{m,n})\le m_{\lambda,d}(D_n)\) for each irreducible
\(L_\lambda\). A strict reverse inequality certifies noncontainment.
B\"urgisser, Ikenmeyer, and Panova show that the occurrence supports cannot
separate these closures in the padded regime \(n\ge m^{25}\) studied in
their theorem~\cite{gctnooccurrence}. A multiplicity comparison can use
unequal positive multiplicities. D\"orfler, Ikenmeyer, and Panova exhibit
such a distinction for a separate orbit-closure problem~\cite{gctmultiplicity}.

\paragraph{Finite symmetry and a polynomial separator.}
On a finite permutation set \(S\), the commutant on \(\mathbb Q^S\)
consists of matrices constant on ordered-pair orbits. Their indicators
give a basis, reused here from Galois Obstruction Theory. The graded
\(G\)-representations above require their own multiplicity calculation.

For the left-right orbit of the \(2\times2\) matrix \(E_{00}\), every
\(LE_{00}R\) is an outer product. The determinant vanishes on the closure
and takes the value one at the identity. Lean checks both the general
separation criterion and this example.

\paragraph{Permanent level spaces.}
The PEqualNP desk asks for a linear bound on the dimension of a linear
space \(L\) with \(\operatorname{per}(J_n+B)=n!\) for every \(B\in L\),
where \(J_n\) is the all-ones matrix. Its finite symmetric test asks for
dimension three at \(n=4\)~\cite{gctlevelbound,gctfour}.
The checked implication \texttt{L-gct-077} converts a dimension bound
\(D\) at \(J_n\) into \(N+D\ge n^2\) for a size-\(N\) border determinantal
expression of \(\operatorname{per}_{n+1}\)~\cite{gctlevelbridge}.
For an anchor \(A\), represent a proposed \(r\)-plane by an independent
frame \(B=(B_1,\ldots,B_r)\). Define \(C_A(B)\) to be the vector of
nonconstant coefficients of
\begin{equation*}
  \operatorname{per}\!\left(A+\sum_{i=1}^r t_i B_i\right).
\end{equation*}
The frame condition is open, given by the rank of the matrix with
columns \(\operatorname{vec}(B_i)\). A constant level is the fiber cut
out by \(C_A(B)=0\). The coefficient comparison gives the pullback
\begin{equation*}
\begin{CD}
 C_A^{-1}(0) @>>> \operatorname{Fr}_r(V) \\
 @VVV @VV{C_A}V \\
 \{0\} @>>> \mathbb A^M,
\end{CD}
\end{equation*}
where \(M\) counts the monomials in the parameters of total degree
between one and \(n\). Every coefficient is required. Keeping only the
linear terms gives the tangent condition.

For \(n\ge 2\), take a symmetric perturbation supported on the first row and column.
Write \(d\) for the perturbation at \((0,0)\) and \(u_j\) for the
mirrored off-diagonal entries, with \(s=\sum_{j>0}u_j\). Splitting the
permanent into permutations that fix index zero and those that do not
gives
\begin{equation*}
  \operatorname{per}(J_n+tB)
   =(n-1)!(1+td)+(n-2)!(n-1+ts)^2.
\end{equation*}
Over \(\mathbb C\), the quadratic coefficient forces \(s=0\), and the
linear coefficient then forces \(d=0\). The constant directions within
this support therefore form an \((n-2)\)-plane. At \(n=4\), the choice
\(u=(a,b,-a-b)\) gives a symmetric two-plane in the fiber at \(24\).
The three-dimensional desk problem requires directions beyond this support.

There is also an upper bound on the entire symmetric fiber. For the
Frobenius bilinear pairing \(\langle B,C\rangle=\sum_{i,j}B_{ij}C_{ij}\),
put \(K=J_n-I_n\). The Hessian is
\begin{equation*}
  \mathrm d^2\!\operatorname{per}_{J_n}(B,C)
    =(n-2)!\,\langle B,KCK\rangle.
\end{equation*}
Two selected entries require distinct rows and columns and have
\((n-2)!\) completions. The first derivative is proportional to the
entry sum, so the symmetric tangent space is \(J_n^\perp\). The Frobenius
pairing on it is nondegenerate since \(\langle J_n,J_n\rangle=n^2\).
The self-adjoint operator \(C\mapsto KCK\) is invertible and preserves
this tangent space by \(KJ_nK=(n-1)^2J_n\). Hence the restricted Hessian
is nondegenerate. A linear level space satisfies \(L\subseteq L^\perp\)
for this Hessian. At \(n=4\) the tangent dimension is nine, so together
with the explicit two-plane we have \(2\le \max\dim L\le 4\).
